% Corrected French for printed p.147.
% Corrections: du truc; il y a donc; induit; K as the argument of R^n(F∘H^0).
\SourcePage{147}

\LayerTarget{tohoku:construction:hyperhomology-resolution:a:p2}
\noindent un objet injectif.
\LayerTarget{tohoku:construction:hyperhomology-resolution:b}
b) On prend une ``résolution injective'' de $K$ au sens de [6, Chap. XVII]
(les guillemets sont nécessaires, la terminologie de [6] étant visiblement
équivoque), i. e. on suppose que les $L^{p,q}$ sont injectifs, et que pour $p$
fixé, si on prend les cycles (resp. bords, resp. la cohomologie) des $L^{p,*}$
par rapport au premier opérateur différentiel, on trouve une résolution
injective du truc formé des cycles (resp. les bords, resp. la cohomologie) dans
$K^p$.
\LayerTarget{tohoku:construction:hyperhomology-identification}
--- Ceci dit, si $L$ est comme ci-dessus, alors $H^*F(L)$ est essentiellement
indépendant du choix de $L$, et s'identifie d'ailleurs à
$R^*(F\mathbin{\circ}H^0)(K)$ (où $H^0$ est considéré comme un foncteur
covariant exact à gauche $\mathfrak K\to \mathbf{C}$, et
$F\mathbin{\circ}H^0$ est le foncteur composé
$\mathfrak K\to \mathbf{C}'$). Pour le voir, il suffit de prendre une
résolution injective $L'$ de $K$ (au sens de a)), il y a donc un homomorphisme
(unique à homotopie près) de la résolution $L$ dans la résolution $L'$, d'où
un homomorphisme $F(L)\to F(L')$, qui induit un \textit{isomorphisme}
$I_2^{p,q}(F(L))\to
I_2^{p,q}(F(L'))$
(les deux membres s'identifient en effet à $H^p(R^qF(K))$), donc un
\textit{isomorphisme} de $HF(L)$ sur $HF(L')$. Or ce dernier s'explicite grâce
à la seconde suite spectrale du bicomplexe $F(L')$ : on a
$H_{I}^q(F(L'))=0$ pour $q>0$ comme on voit aussitôt, on trouve donc
$H^nF(L')=H^n(FH^0(L'^{*n}))$,
or le deuxième membre est par définition $R^n(F\mathbin{\circ}H^0)(K)$.
D'où la définition et le mode de calcul général de \textit{l'hyperhomologie}
$\mathfrak R F(K)=R^*(FH^0)(K)$ du foncteur $F$ par rapport au
complexe $K$, et de la première suite spectrale, de terme initial
$I_2^{p,q}F(K)=H^p(R^qF(K))$ qui y aboutit.
Si maintenant on prend pour $L$ une ``résolution injective'' de $K$ comme dans
b), alors la deuxième suite spectrale du bicomplexe $F(L)$ est essentiellement
indépendante de $L$ (puisque $L$ est unique à une équivalence d'homotopie
près), elle aboutit à $\mathfrak R F(K)$ et son terme initial est celui donné
dans (2.4.2) ; il suffirait d'ailleurs qu'on ait $R^nF(L^{pq})=0$ pour $n>0$
(au lieu de $L^{p,q}$ injectif) dans l'énoncé des conditions b), pour que la
deuxième suite spectrale du complexe $F(L)$ soit celle envisagée dans (2.4.2).
\LayerTarget{tohoku:remark:derived-spectral-functor-scope}
Notons encore que, si les degrés de $K$ sont positifs, les deux suites
spectrales dérivées de $F$ sont des suites spectrales \textit{cohomologiques}.
On peut encore définir les suites spectrales dérivées de $F$ sur $K$ si on ne
suppose plus que $K$ est à degrés bornés inférieurement, \textit{pourvu que
$F$ soit de dimension injective finie}, i. e. $R^pF=0$ pour $p$ convenable. Ce
fait ne semble pas contenu dans [6], mais comme nous ne nous en servirons pas
dans la suite, nous nous bornons à le signaler ici sans démonstration. --- Bien
entendu, on pourrait définir aussi des foncteurs spectraux dérivés à gauche de
$F$, pourvu que $\mathbf{C}$ contienne assez d'objets projectifs, et on peut envisager
aussi le cas d'un foncteur contravariant. Dans [6], on envisage le cas où $F$
est un multifoncteur, mais nous n'aurons pas à nous en servir.

\LayerTarget{tohoku:construction:composite-functor-spectral-sequences:p1}
Soient $\mathbf{C},\mathbf{C}',\mathbf{C}''$ trois catégories abéliennes, considérons deux foncteurs
covariants $F:\mathbf{C}\to \mathbf{C}'$ et $G:\mathbf{C}'\to \mathbf{C}''$. On suppose
que tout objet de $\mathbf{C}$ ou $\mathbf{C}'$ est isomorphe à un sous-truc d'un objet injectif,
ce qui permet de considérer les foncteurs dérivés à droite de $F,G,GF$, on se
propose d'établir des relations entre ces foncteurs dérivés. Soit $A\in \mathbf{C}$,
soit $C(A)$ le complexe associé à une résolution droite de $A$ par des objets
injectifs, considérons le complexe $F(C(A))$ dans $\mathbf{C}'$, il est déterminé à une
homotopie près quand on fait varier

\TohokuContinuation{148}
